\documentclass[11pt,a4paper]{article}
\usepackage[T1]{fontenc}
\usepackage[utf8]{inputenc}
\usepackage[ngerman]{babel}
\usepackage[margin=2.5cm]{geometry}
\usepackage[hidelinks]{hyperref}
\usepackage{array}
\setlength{\parindent}{0pt}
\setlength{\emergencystretch}{2em}
\renewcommand{\arraystretch}{1.25}
\sloppy
\begin{document}
\shorthandoff{"}
\begin{center}
  {\Large\textbf{Webtechnologien -- Übungsblatt 5}}
\end{center}
\section*{Aufgabe 5: Retrieving Products via the Fetch API}
\textbf{Aufgabe:} Wie Aufgabe 4, aber die Produkte mit der Fetch API
laden und Fehler richtig behandeln.

\medskip
\textbf{Ergebnis:} Gleicher Aufbau wie in Aufgabe 4: \texttt{products.js}
ist die Originaldatei mit \texttt{[5 GEÄNDERT]} (URL),
\texttt{[5 VERSCHOBEN]} (Aufruf) und \texttt{[5 ENTFERNT]} (leere
Vorlage); alles Neue steht nach fünf Leerzeilen am Dateiende. Neu ist nur
\texttt{initialize}, \texttt{renderProducts}, \texttt{showError} und der
\texttt{<style>}-Block in \texttt{products.html} sind gleich wie in 4.
\begin{itemize}
  \item \texttt{fetch(jsonURL)} schickt die Anfrage ab und gibt sofort
    ein Promise zurück.
  \item Erstes \texttt{.then(response => ...)}: \texttt{response} ist die
    Antwort des Servers. \textbf{Wichtig:} \texttt{fetch} lehnt bei
    404 oder 500 \emph{nicht} von selbst ab, das Promise wird trotzdem
    erfüllt. Deshalb wird \texttt{response.ok} (Status 200--299) geprüft
    und sonst selbst ein Fehler geworfen
    (\texttt{throw new Error("HTTP 404 Not Found")}). Ist alles gut,
    liest \texttt{return response.json()} den Inhalt als JSON; das ist
    wieder ein Promise.
  \item Zweites \texttt{.then(products => renderProducts(...))}: Jetzt
    ist das fertige Array da und wird wie in 4 angezeigt.
  \item \texttt{.catch(error => showError(...))}: \emph{ein} catch am
    Ende fängt alles ab, was davor schiefgeht: der selbst geworfene
    HTTP-Fehler, kaputtes JSON aus \texttt{response.json()} und eine
    Netzwerkfehler-Ablehnung von \texttt{fetch} (Server aus,
    \texttt{file://}). Der Benutzer sieht eine rote Meldung, die Seite
    stürzt nicht ab.
\end{itemize}

\textbf{Unterschied zu Aufgabe 4:} \texttt{XMLHttpRequest} arbeitet mit
Events (\texttt{onload}, \texttt{onerror}) und braucht für
\texttt{JSON.parse} ein eigenes \texttt{try}/\texttt{catch}.
\texttt{fetch} arbeitet mit Promises wie in Aufgabe 3, alle Fehler landen
im selben \texttt{.catch}. Bei beiden muss man den HTTP-Status selbst
prüfen.

\medskip
\textbf{Hinweis:} In XAMPP liegt der Ordner unter
\texttt{C:\textbackslash xampp\textbackslash htdocs\textbackslash ex5},
aufrufen mit \url{http://localhost/ex5/products.html}. Wegen der
relativen URL \texttt{"products.json"} funktioniert es in jedem Ordner.

\medskip
\textbf{Test:} In Chrome:
\begin{center}
\begin{tabular}{|>{\raggedright\arraybackslash}p{5.2cm}|>{\raggedright\arraybackslash}p{8.5cm}|}
\hline
\textbf{Fall} & \textbf{Ergebnis} \\ \hline
\url{http://localhost/ex5/products.html} (XAMPP) & 4 Produkte, gleich
  wie in Aufgabe 4 \\ \hline
Datei nicht gefunden (404) & \glqq Products could not be loaded: HTTP 404
  Not Found\grqq \\ \hline
kaputtes JSON & Meldung mit dem JSON-Fehler \\ \hline
Server nicht erreichbar & Meldung \glqq Failed to fetch\grqq \\ \hline
über \texttt{file://} geöffnet & Meldung \glqq Failed to fetch\grqq{} \\ \hline
\end{tabular}
\end{center}
Kein unbehandelter Fehler: In der Konsole stehen nur die Meldung des
Browsers (z.\,B. \glqq Failed to load resource\grqq{} bzw. bei
\texttt{file://} die CORS-Sperre) und die eigene Meldung von
\texttt{showError}.

\medskip
{\small \textbf{Hilfsmittel:} Claude Code (Claude Opus 5.5) für Erstellung, Prüfung und Dokumentation.}
\end{document}
